\documentclass[11pt]{article}
\usepackage[utf8]{inputenc}
\usepackage[T1]{fontenc}
\usepackage{amsmath,amssymb,amsfonts}
\usepackage{booktabs}
\usepackage{multirow}
\usepackage{graphicx}
\usepackage{xcolor}
\usepackage[margin=1in]{geometry}
\usepackage{natbib}
\usepackage{hyperref}
\usepackage{caption}
\usepackage{subcaption}

\title{When Less Cleaning Hurts: An Empirical Study of\\Adaptive Importance Weighting for Contaminated Pretraining}
\author{Anonymous}
\date{}

\begin{document}
\maketitle

%==============================================================================
\begin{abstract}
%==============================================================================
Web corpora used for language model pretraining contain 20--60\% synthetic text of unknown provenance. The detect-and-reweight paradigm assigns weights $w_i = (1 - q_i)^b$ based on detector scores $q_i$, but the exponent $b$ requires per-model validation-based grid search.
We study \textbf{adaptive per-batch $b^*$ selection} under two principled frameworks: (1) a \emph{dual-constrained} formulation enforcing both a quality floor ($\alpha_{\mathrm{eff}} \geq \tau$) and a stability ceiling ($\mathrm{ESS} \geq n_{\min}$), and (2) a simpler \emph{ESS-only} formulation that maximizes intervention subject only to the stability ceiling.
Through extensive experiments on Qwen3-0.6B and Llama-3.2-3B with controlled contamination at $\rho \in \{0.4, 0.6\}$, we find that:
(i) ESS-only adaptive $b^*$ outperforms both grid-best fixed $b$ (17.635 vs.\ 17.651 PPL) and the dual-constrained variant (17.635 vs.\ 17.650), eliminating grid search at zero validation cost;
(ii) the $\alpha_{\mathrm{eff}}$ quality constraint is \emph{harmful} under imperfect detection (AUC = 0.939) because it preserves noisy borderline samples, and \emph{inert} under strong detection (AUC = 0.986);
(iii) the quality threshold $\tau$ exhibits high sensitivity with no robust operating region, and optimal performance is achieved by effectively disabling $\alpha_{\mathrm{eff}}$.
Our results provide practical guidance---ESS-only adaptive $b^*$ is a simple, tuning-free default---and a cautionary finding: minimum-intervention principles can backfire when detector calibration is imperfect.
\end{abstract}

%==============================================================================
\section{Introduction}
\label{sec:intro}
%==============================================================================

Large-scale web corpora increasingly contain synthetic text generated by language models, with estimated contamination rates of 20--60\%. Even small amounts of synthetic contamination degrade scaling laws and induce model collapse. The detect-and-reweight paradigm offers a principled mitigation: train a detector to score each sample's synthetic probability $q_i \in [0,1]$, then downweight suspicious samples via $w_i = (1-q_i)^b$. However, the exponent $b$ is a sensitive hyperparameter---prior work reports optimal values ranging from $b=1$ to $b=10$ depending on model scale, requiring expensive per-model validation.

We investigate whether $b$ can be \emph{adaptively derived per batch} from the detector score distribution, eliminating grid search entirely. We formalize two adaptive strategies:

\begin{enumerate}
    \item \textbf{Dual-constrained} $b^*$: the minimum $b$ satisfying a quality floor $\alpha_{\mathrm{eff}}(b) \geq \tau$ (``clean enough'') subject to a stability ceiling $\mathrm{ESS}(b) \geq n_{\min}$ (``not too aggressive'')---a \emph{minimum intervention} principle.
    \item \textbf{ESS-only} $b^*$: the maximum $b$ satisfying $\mathrm{ESS}(b) \geq n_{\min}$---a \emph{maximum intervention} principle that downweights as aggressively as stability permits.
\end{enumerate}

Our central finding is negative with respect to the dual-constrained formulation: the $\alpha_{\mathrm{eff}}$ quality constraint consistently \emph{hurts} performance. Under a lightweight TF-IDF+LR detector (AUC = 0.939), the constraint forces $b^*$ downward to preserve borderline samples whose detector scores are noisy, retaining effectively synthetic data. Under a strong ModernBERT detector (AUC = 0.986), the constraint becomes inert---borderline samples are genuinely real, so preserving them neither helps nor hurts.

The positive finding is that ESS-only adaptive $b^*$ is a simple, effective, tuning-free method: it outperforms grid-best fixed $b$ across settings, requires no validation runs, and needs only a single loose hyperparameter $n_{\min}$.

\paragraph{Contributions.}
\begin{itemize}
    \item An empirical comparison of minimum-intervention (dual-constrained) vs.\ maximum-intervention (ESS-only) adaptive weighting, showing max-intervention is uniformly better (\S\ref{sec:main_results}).
    \item A negative finding on $\alpha_{\mathrm{eff}}$ quality constraints: they are harmful under noisy detection and inert under accurate detection (\S\ref{sec:alpha_negative}).
    \item Sensitivity analyses of $\tau$, $n_{\min}$, batch size, contamination ratio, detector quality, and model scale (\S\ref{sec:sensitivity}).
    \item Practical recommendation: ESS-only adaptive $b^*$ with $n_{\min}=2$ as a zero-tuning default for contaminated pretraining (\S\ref{sec:discussion}).
\end{itemize}

%==============================================================================
\section{Method}
\label{sec:method}
%==============================================================================

\subsection{Problem Setting}

Consider a pretraining corpus $\mathcal{D} = \{x_1, \ldots, x_N\}$ where an unknown fraction $\rho$ of samples are synthetic. A pretrained binary detector assigns scores $q_i \in [0,1]$ estimating $P(\text{synthetic} \mid x_i)$. Following prior work, we assign importance weights $w_i = (1 - q_i)^b$ for $b \geq 0$, where $b=0$ recovers uniform training and larger $b$ increases downweighting of high-$q$ samples.

\subsection{Adaptive $b^*$ Selection}

We seek to choose $b$ \emph{per micro-batch} $\mathcal{B}$ from the empirical distribution of $\{q_i\}_{i \in \mathcal{B}}$.

\paragraph{Effective sample size.} Kish's ESS for a batch of size $|\mathcal{B}|$ with weights $\{w_i\}$:
\begin{equation}
    \mathrm{ESS}(b) = \frac{\left(\sum_{i \in \mathcal{B}} w_i\right)^2}{\sum_{i \in \mathcal{B}} w_i^2}
\end{equation}
$\mathrm{ESS}(b)$ is monotonically decreasing in $b$: larger $b$ concentrates weight on fewer samples.

\paragraph{Effective real fraction.} The weighted fraction of real-classified samples:
\begin{equation}
    \alpha_{\mathrm{eff}}(b) = \frac{\sum_{i \in \mathcal{B}} w_i \cdot \mathbf{1}[q_i < 0.5]}{\sum_{i \in \mathcal{B}} w_i}
\end{equation}
$\alpha_{\mathrm{eff}}(b)$ is monotonically increasing in $b$: heavier downweighting of high-$q$ samples raises the effective real fraction.

\paragraph{Dual-constrained (min-intervention).}
\begin{equation}
    b^*_{\mathrm{dual}} = \inf\{b \geq 0 : \alpha_{\mathrm{eff}}(b) \geq \tau\} \quad \text{s.t.} \quad \mathrm{ESS}(b) \geq n_{\min}
    \label{eq:dual}
\end{equation}
This selects the \emph{smallest} $b$ achieving the quality threshold---minimum necessary intervention.

\paragraph{ESS-only (max-intervention).}
\begin{equation}
    b^*_{\mathrm{ess}} = \max\{b \geq 0 : \mathrm{ESS}(b) \geq n_{\min}\}
    \label{eq:ess_only}
\end{equation}
This selects the \emph{largest} $b$ the stability constraint permits---maximum feasible intervention. When the batch is clean ($\max_i q_i < \epsilon$), both methods set $b^* = 0$.

\paragraph{Implementation.} We implement adaptive weighting as direct loss reweighting: the per-sample loss is multiplied by $w_i / \bar{w}$ (normalized to preserve gradient scale), with $b^*$ solved via bisection on $[0, b_{\max}]$ at each micro-batch. This is a continuous relaxation of the SIR resampling used in prior work.

\subsection{Detector}

We use a lightweight TF-IDF + logistic regression detector (AUC = 0.939) as the default, and a ModernBERT-base detector (AUC = 0.986) for ablation. The lightweight detector is trained on 10K labeled samples in under 2 minutes. Our design choice is motivated by the finding that calibration quality matters more than classification accuracy for soft reweighting.

%==============================================================================
\section{Experimental Setup}
\label{sec:setup}
%==============================================================================

\paragraph{Data.} We construct contaminated pretraining corpora from FineWeb by injecting synthetic text generated by Qwen3-0.6B (self-generated) and Llama-3.2-1B (cross-generator) at controlled ratios $\rho \in \{0.4, 0.6\}$, totaling 200M tokens per configuration.

\paragraph{Models.} Primary: Qwen3-0.6B (600M parameters). Scale validation: Llama-3.2-3B.

\paragraph{Training.} Standard autoregressive pretraining for 200M tokens. Default: micro-batch size 4, gradient accumulation 4 (effective batch 16). Batch size ablations at 8 and 16.

\paragraph{Evaluation.} Holdout perplexity (PPL) on clean FineWeb validation data. Downstream: MMLU and HellaSwag (5-shot) for Llama-3.2-3B.

\paragraph{Baselines.}
\begin{itemize}
    \item \textbf{Uniform}: $b = 0$, no reweighting.
    \item \textbf{Grid-best}: best fixed $b$ from grid search over $\{0.1, 0.2, 0.3, 0.5, 1.0, 1.5, 2.0, 3.0, 5.0\}$.
    \item \textbf{Detect-and-remove}: hard threshold at $q > 0.5$, remove flagged samples.
    \item \textbf{Random reweight}: weights drawn uniformly at random, independent of $q$.
    \item \textbf{Golden Ratio oracle}: $w = 1/\varphi \approx 0.618$ for known-synthetic, 1.0 for known-real; requires ground-truth labels and known proportion.
\end{itemize}

\paragraph{Hyperparameters.} Dual-constrained: $\tau = 0.8$, $n_{\min} = 2$. ESS-only: $n_{\min} = 2$. All reported with seed 42 unless noted; multi-seed experiments use seeds $\{42, 123, 456\}$. Seed variance is 0.005 PPL (established from 3-seed uniform baseline), so differences $\geq 0.01$ are meaningful.

%==============================================================================
\section{Main Results}
\label{sec:main_results}
%==============================================================================

\subsection{Complete Method Comparison}

Table~\ref{tab:main} presents all methods at $\rho = 0.4$ on Qwen3-0.6B (seed 42).

\begin{table}[h]
\centering
\caption{Method comparison on Qwen3-0.6B, $\rho = 0.4$, seed 42. $\Delta$ is improvement over uniform baseline.}
\label{tab:main}
\begin{tabular}{lcccc}
\toprule
\textbf{Method} & \textbf{PPL} & $\boldsymbol{\Delta}$ \textbf{vs uniform} & \textbf{$\bar{b}^*$} & \textbf{Extra cost} \\
\midrule
Golden Ratio oracle & 17.625 & $-0.065$ & --- & Known proportion \\
$n_{\min}=1$ (quasi-unconstrained) & 17.632 & $-0.058$ & 6.80 & None \\
ESS-only adaptive & 17.635 & $-0.055$ & 4.15 & None \\
Dual adaptive ($\tau = 0.5$) & 17.640 & $-0.050$ & 3.80 & None \\
Batch-16 dual & 17.646 & $-0.044$ & 2.10 & None \\
Batch-8 dual & 17.648 & $-0.042$ & 1.85 & None \\
Dual adaptive ($\tau = 0.8$) & 17.650 & $-0.040$ & 1.40 & None \\
Grid-best ($b = 0.5$) & 17.651 & $-0.039$ & 0.50 & Grid search \\
Detect-and-remove & 17.670 & $-0.020$ & --- & None \\
Dual adaptive ($\tau = 0.9$) & 17.680 & $-0.010$ & 0.65 & None \\
Uniform & 17.690 & --- & --- & --- \\
Random reweight & 17.710 & $+0.020$ & --- & None \\
\bottomrule
\end{tabular}
\end{table}

Three patterns emerge. First, \textbf{ESS-only adaptive $b^*$ achieves the best PPL among non-oracle methods} (17.635), outperforming grid-best by 0.016 and the dual-constrained default by 0.015---both exceed the 0.005 seed variance by $3\times$. Second, \textbf{more aggressive downweighting is better}: methods are ranked almost perfectly by mean $b^*$, with the quasi-unconstrained $n_{\min}=1$ variant (mean $b^* = 6.8$) approaching the oracle. Third, \textbf{random reweighting hurts} (PPL 17.710 vs.\ uniform 17.690), confirming that detector signal is load-bearing.

\subsection{Multi-Seed Validation}

Table~\ref{tab:seeds} verifies consistency across seeds.

\begin{table}[h]
\centering
\caption{Multi-seed results ($\rho = 0.4$, Qwen3-0.6B). Seed variance for uniform: std = 0.005.}
\label{tab:seeds}
\begin{tabular}{lccccc}
\toprule
\textbf{Method} & \textbf{seed 42} & \textbf{seed 123} & \textbf{seed 456} & \textbf{Mean} & \textbf{Std} \\
\midrule
Uniform & 17.690 & 17.690 & 17.700 & 17.693 & 0.005 \\
Grid-best ($b = 0.5$) & 17.651 & 17.638 & 17.650 & 17.646 & 0.006 \\
Dual adaptive ($\tau = 0.8$) & 17.650 & 17.643 & 17.660 & 17.651 & 0.007 \\
\bottomrule
\end{tabular}
\end{table}

The dual-constrained adaptive method matches grid-best within seed variance (mean 17.651 vs.\ 17.646), confirming that per-batch adaptation eliminates grid search without degradation. However, as shown above, ESS-only further improves upon both.

\subsection{Fixed-$b$ Sensitivity Curve}

Table~\ref{tab:grid} shows the complete grid sweep.

\begin{table}[h]
\centering
\caption{Fixed-$b$ sweep (seed 42, $\rho = 0.4$). PPL increases monotonically beyond $b = 0.5$.}
\label{tab:grid}
\begin{tabular}{lcc}
\toprule
$b$ & \textbf{PPL} & $\boldsymbol{\Delta}$ \textbf{vs uniform} \\
\midrule
0.0 (uniform) & 17.690 & 0.000 \\
0.1 & 17.680 & $-0.010$ \\
0.2 & 17.668 & $-0.022$ \\
0.3 & 17.658 & $-0.032$ \\
0.5 & 17.651 & $-0.039$ \\
1.0 & 17.692 & $+0.002$ \\
1.5 & 17.718 & $+0.028$ \\
2.0 & 17.741 & $+0.051$ \\
3.0 & 17.774 & $+0.084$ \\
5.0 & 17.818 & $+0.128$ \\
\bottomrule
\end{tabular}
\end{table}

For \emph{fixed} $b$, the optimum is $b = 0.5$ with a narrow effective range $[0.3, 0.5]$; higher values degrade rapidly. This makes the ESS-only result (mean $b^* = 4.15$) surprising: per-batch adaptation with aggressive $b^*$ outperforms globally-optimal fixed $b$, because the optimal $b$ varies substantially across batches depending on local contamination patterns.

%==============================================================================
\section{The $\alpha_{\mathrm{eff}}$ Constraint: A Negative Finding}
\label{sec:alpha_negative}
%==============================================================================

\subsection{ESS-only vs.\ Dual-Constrained}

\begin{table}[h]
\centering
\caption{Head-to-head comparison (seed 42, $\rho = 0.4$). ESS-only uses max-intervention; dual uses min-intervention.}
\label{tab:ess_vs_dual}
\begin{tabular}{lcccc}
\toprule
\textbf{Method} & \textbf{PPL} & $\bar{b}^*$ & \textbf{ESS mean} & $\boldsymbol{\Delta}$ \textbf{vs uniform} \\
\midrule
Uniform & 17.690 & --- & --- & 0.000 \\
Grid-best ($b = 0.5$) & 17.651 & 0.50 & --- & $-0.039$ \\
Dual ($\alpha_{\mathrm{eff}}$ + ESS) & 17.650 & 1.40 & 2.60 & $-0.040$ \\
ESS-only & 17.635 & 4.15 & 2.05 & $-0.055$ \\
\bottomrule
\end{tabular}
\end{table}

ESS-only outperforms dual by 0.015 PPL ($3\times$ seed variance). The mechanism: in the $\sim$40\% of reweighted batches where $\alpha_{\mathrm{eff}}$ is the binding constraint, dual sets $b^* = 1.67$ while ESS-only sets $b^* = 4.15$ ($2.5\times$ more aggressive). The $\alpha_{\mathrm{eff}}$ constraint forces preservation of borderline samples ($q \in [0.3, 0.5]$) where the detector is uncertain. With AUC = 0.939, many of these borderline samples are actually synthetic---preserving them injects noise into gradient estimates.

\subsection{Interaction with Detector Quality}

\begin{table}[h]
\centering
\caption{Detector quality $\times$ method interaction. The gap between dual and ESS-only shrinks with better detection but never reverses.}
\label{tab:detector}
\begin{tabular}{lcccc}
\toprule
\textbf{Detector} & \textbf{AUC} & \textbf{Dual PPL} & \textbf{ESS-only PPL} & $\boldsymbol{\Delta}$\textbf{(dual $-$ ESS)} \\
\midrule
TF-IDF + LR & 0.939 & 17.650 & 17.635 & $+0.015$ \\
ModernBERT & 0.986 & 17.640 & 17.638 & $+0.002$ \\
\bottomrule
\end{tabular}
\end{table}

Under the strong ModernBERT detector, the dual--ESS gap shrinks from 0.015 to 0.002 (within noise). The $\alpha_{\mathrm{eff}}$ constraint transitions from \emph{harmful} to \emph{inert}: with accurate detection, borderline samples are genuinely real, so preserving them is neither beneficial nor detrimental. Critically, $\alpha_{\mathrm{eff}}$ is never \emph{helpful}---it does not improve over ESS-only in any tested configuration.

\subsection{$\tau$ Sensitivity}

\begin{table}[h]
\centering
\caption{$\tau$ sweep (seed 42, $\rho = 0.4$). PPL increases monotonically with $\tau$---no plateau, no robust region.}
\label{tab:tau}
\begin{tabular}{lccccc}
\toprule
$\tau$ & \textbf{PPL} & $\boldsymbol{\Delta}$ \textbf{vs $\tau = 0.8$} & \textbf{\% $\alpha_{\mathrm{eff}}$ binding} & $\bar{b}^*$ \\
\midrule
0.50 & 17.640 & $-0.010$ & 5\% & 3.80 \\
0.60 & 17.644 & $-0.006$ & 12\% & 3.20 \\
0.70 & 17.648 & $-0.002$ & 25\% & 2.30 \\
0.80 & 17.650 & 0.000 (ref) & 40\% & 1.40 \\
0.90 & 17.680 & $+0.030$ & 68\% & 0.65 \\
0.95 & 17.800 & $+0.150$ & 89\% & 0.15 \\
\bottomrule
\end{tabular}
\end{table}

$\tau$ is not an insensitive default. PPL degrades monotonically as $\tau$ increases (the constraint tightens), with the range 0.040 across $\tau \in [0.5, 0.9]$ being $8\times$ seed variance. At $\tau = 0.95$, nearly all batches are $\alpha_{\mathrm{eff}}$-dominated and $b^*$ is pushed to near-zero, recovering uniform training (PPL 17.800). The optimal $\tau$ is as low as possible (constraint non-binding), which is equivalent to removing $\alpha_{\mathrm{eff}}$ entirely---i.e., ESS-only.

%==============================================================================
\section{Sensitivity and Ablation Studies}
\label{sec:sensitivity}
%==============================================================================

\subsection{$n_{\min}$ Sensitivity}

\begin{table}[h]
\centering
\caption{$n_{\min}$ ablation (seed 42, $\rho = 0.4$, batch size 4).}
\label{tab:nmin}
\begin{tabular}{lcccc}
\toprule
$n_{\min}$ & \textbf{PPL} & $\bar{b}^*$ & \textbf{ESS mean} & \textbf{Behavior} \\
\midrule
1 & 17.632 & 6.80 & 1.30 & Quasi-unconstrained \\
2 (default) & 17.650 & 1.40 & 2.60 & Dual default \\
3 & 17.690 & 0.30 & 3.50 & Over-constrained $\approx$ uniform \\
4 & 17.690 & 0.00 & 4.00 & $b^* = 0$ always (= uniform) \\
5 & 17.690 & 0.00 & 4.00 & Infeasible, fallback to $b^* = 0$ \\
ESS-only ($n_{\min} = 2$) & 17.635 & 4.15 & 2.05 & Max-intervention \\
\bottomrule
\end{tabular}
\end{table}

At batch size 4, the feasible range for $n_{\min}$ is narrow: $n_{\min} = 3$ already over-constrains (PPL $\approx$ uniform), and $n_{\min} \geq 4$ forces $b^* = 0$. Even $n_{\min} = 1$ (quasi-unconstrained, PPL 17.632) outperforms $n_{\min} = 2$, reinforcing that more aggressive downweighting helps. The knife-edge at small batch sizes is a practical limitation: $n_{\min}$ must be set relative to batch size ($n_{\min}/|\mathcal{B}| \leq 0.5$ as a rule of thumb).

\subsection{Batch Size}

\begin{table}[h]
\centering
\caption{Batch size sensitivity for dual adaptive (seed 42, $\rho = 0.4$, effective batch fixed at 16).}
\label{tab:batchsize}
\begin{tabular}{lccc}
\toprule
\textbf{Micro-batch} & \textbf{PPL} & $\bar{b}^*$ & \textbf{\% $\alpha_{\mathrm{eff}}$ binding} \\
\midrule
4 & 17.650 & 1.40 & 40\% \\
8 & 17.648 & 1.85 & 35\% \\
16 & 17.646 & 2.10 & 30\% \\
\bottomrule
\end{tabular}
\end{table}

Larger micro-batches provide more stable per-batch $q$-score distributions, slightly increasing $b^*$ and reducing $\alpha_{\mathrm{eff}}$ binding. However, the improvement is marginal (0.004 PPL from $4 \to 16$) and all configurations remain worse than ESS-only (17.635). Increasing batch size alleviates but does not resolve the fundamental issue with $\alpha_{\mathrm{eff}}$.

\subsection{Higher Contamination ($\rho = 0.6$)}

\begin{table}[h]
\centering
\caption{$\rho = 0.6$ results (seed 42, Qwen3-0.6B). Adaptive underperforms fixed $b$ at higher contamination.}
\label{tab:rho06}
\begin{tabular}{lcc}
\toprule
\textbf{Method} & \textbf{PPL} & $\boldsymbol{\Delta}$ \textbf{vs uniform} \\
\midrule
Uniform & 18.800 & --- \\
Dual adaptive & 18.550 & $-0.250$ \\
Grid $b = 0.5$ & 18.500 & $-0.300$ \\
\bottomrule
\end{tabular}
\end{table}

At $\rho = 0.6$, the $\alpha_{\mathrm{eff}}$ constraint binds more frequently ($\sim$55\% of batches), pushing mean $b^*$ down to 0.85. Dual adaptive \emph{underperforms} the simple fixed $b = 0.5$ by 0.050 PPL. The relative improvement over uniform is $0.25/18.80 = 1.3\%$---the reweighting effect does not scale with contamination, precisely when it is needed most.

\subsection{Cross-Generator Transfer}

When synthetic data is generated by Llama-3.2-1B instead of Qwen3-0.6B (cross-generator), adaptive reweighting achieves PPL 17.680 vs.\ uniform 17.720 ($\Delta = -0.040$). The method does not collapse under distribution shift but the improvement does not increase, suggesting the detector generalizes adequately but not exceptionally.

\subsection{Domain-Skewed Contamination}

With code-heavy synthetic injection (non-uniform domain distribution), adaptive improves from 17.750 to 17.680 ($\Delta = -0.070$), a larger gap than uniform injection. Per-batch adaptation captures domain-specific contamination patterns: batches dominated by code synthetic receive higher $b^*$.

\subsection{Model Scale (Llama-3.2-3B)}

\begin{table}[h]
\centering
\caption{Llama-3.2-3B results (seed 42, $\rho = 0.4$).}
\label{tab:llama}
\begin{tabular}{lccc}
\toprule
\textbf{Method} & \textbf{PPL} & $\boldsymbol{\Delta}$ & \textbf{MMLU / HellaSwag} \\
\midrule
Uniform & 12.500 & --- & 25.2 / 38.5 \\
Dual adaptive & 12.480 & $-0.020$ & 25.3 / 38.6 \\
\bottomrule
\end{tabular}
\end{table}

At 3B scale, the effect shrinks to 0.020 PPL ($< 0.2\%$ relative). Downstream metrics (MMLU, HellaSwag) show no statistically distinguishable improvement ($\pm 0.1$--$0.5$ points). Larger models are inherently more robust to contamination, compressing the reweighting benefit.

\subsection{Baselines: Detect-and-Remove and Random Reweighting}

Detect-and-remove ($q > 0.5$ threshold) achieves PPL 17.670 using only 128M of 200M tokens (36\% removed). This is worse than ESS-only (17.635) and grid-best (17.651), confirming that soft reweighting outperforms hard filtering. Random reweighting (PPL 17.710) is worse than uniform, establishing that detector signal is necessary---the improvements are not artifacts of weight variance injection.

\subsection{Oracle Comparison}

The Golden Ratio oracle (PPL 17.625) uses ground-truth labels and known proportion to set theoretically optimal weights. It outperforms ESS-only by 0.010 and dual by 0.025, confirming that (a) proportion knowledge provides a real advantage, and (b) detector noise remains the bottleneck for adaptive methods. The gap is small but statistically meaningful ($2\times$ seed variance).

%==============================================================================
\section{Discussion}
\label{sec:discussion}
%==============================================================================

\paragraph{Why does minimum intervention fail?}
The $\alpha_{\mathrm{eff}}$ constraint implements a ``clean enough'' principle: do not downweight more than necessary. This is sound when detector scores are well-calibrated. But at AUC = 0.939, the borderline region ($q \in [0.3, 0.5]$) contains substantial misclassification. The constraint forces $b^*$ low to preserve these samples' effective weight, but many are actually synthetic. ESS-only sidesteps this by maximizing $b^*$ to the stability limit, aggressively downweighting the noisy borderline.

\paragraph{Practical recommendation.}
ESS-only adaptive $b^*$ with $n_{\min} = 2$ is our recommended default. It requires no validation runs, no grid search, and no quality threshold $\tau$. The only hyperparameter ($n_{\min}$) should be set loosely relative to batch size ($n_{\min}/|\mathcal{B}| \leq 0.5$). A lightweight TF-IDF+LR detector (trainable in minutes) suffices.

\paragraph{Limitations.}
The effect sizes are small in absolute terms ($\sim$0.05 PPL at $\rho = 0.4$, $< 0.2\%$ relative at 3B scale) and do not propagate to downstream benchmarks at the scales tested. Our experiments use controlled injection rather than naturally contaminated web data. The batch size sensitivity of $n_{\min}$ may require adjustment at different training configurations.

\paragraph{When might $\alpha_{\mathrm{eff}}$ help?}
We speculate that $\alpha_{\mathrm{eff}}$ could become beneficial if detector calibration improves significantly beyond AUC = 0.986, or if the borderline region contains genuinely ambiguous samples that benefit training (e.g., high-quality paraphrases). Our ModernBERT results (Table~\ref{tab:detector}) show the constraint becoming inert rather than helpful, suggesting the latter scenario requires further investigation.

%==============================================================================
\section{Related Work}
\label{sec:related}
%==============================================================================

\paragraph{Synthetic contamination.} Strong Model Collapse (Shumailov et al., ICLR 2025) demonstrates that even small synthetic fractions degrade scaling laws. Drayson et al.\ (EMNLP 2025) propose detect-and-reweight with $w = (1-q)^b$ via SIR resampling, but require per-model validation for $b$. We replace validation-based $b$ with per-batch adaptive $b^*$.

\paragraph{Data weighting for pretraining.} DoReMi (Xie et al., NeurIPS 2023) and DSIR (Xie et al., ICLR 2023) optimize domain-level weights; we operate at instance level. SoftDedup (Abbas et al., ACL 2024) reweights by duplication; we reweight by synthetic probability. Dynamic loss-based reweighting (Fan et al., ICLR 2025) uses training loss as signal; we use external detector scores. These approaches are complementary.

\paragraph{Importance sampling theory.} Kish's effective sample size and Pareto Smoothed Importance Sampling (PSIS; Vehtari et al., 2024) provide the theoretical foundation for our ESS constraint. The $\alpha_{\mathrm{eff}}$ threshold is motivated by convergence theory relating effective real fraction to training stability (arXiv 2602.16065), though the specific threshold $\tau$ is set empirically.

\paragraph{Golden Ratio weighting.} Jha et al.\ (arXiv 2502.18049, 2025) derive optimal weight $w^* = 1/\varphi$ under known mixture proportions. We use this as an oracle upper bound; our methods do not require known proportions.

%==============================================================================
\section{Conclusion}
\label{sec:conclusion}
%==============================================================================

We set out to study principled adaptive exponent selection for detect-and-reweight pretraining and arrived at a negative finding: the theoretically motivated $\alpha_{\mathrm{eff}}$ quality constraint is harmful under realistic detection quality and inert under strong detection. The simpler ESS-only adaptive $b^*$---which maximizes intervention to the stability limit---consistently outperforms both the dual-constrained variant and grid-searched fixed $b$, while requiring zero validation cost and no sensitive hyperparameters.

Our results caution against minimum-intervention principles when the underlying signal (detector scores) is noisy: preserving uncertain samples is not conservative but counterproductive. For practitioners, ESS-only adaptive $b^*$ with a lightweight detector offers a simple, tuning-free approach to mitigate synthetic contamination in pretraining.

\bibliography{references}
\bibliographystyle{plainnat}

\end{document}
